% SEGMENT OLP-0046-S01
% Part: sets-functions-relations
% Chapter: arithmetization
% Section: reflections

\documentclass[../../../include/open-logic-section]{subfiles}

\begin{document}

\olfileid{sfr}{arith}{ref}
\olsection{சில மெய்யியல் சிந்தனைகள்}

தொழில்நுட்ப விவரங்கள் இவ்வளவே. ஆனால் அவை எதைச் சாதித்தன?

மறுப்பதற்கு இடமில்லாத வகையில், அவை சில {அழகிய} தூய கணிதத்தை
நமக்குத் தந்தன. மேலும், சில ஆழமான கருத்தியல் சாதனைகளும் இருந்தன.
மெய்யெண்களுக்கும் விகிதமுறு எண்களுக்கும் இடையிலான முக்கிய வேறுபாட்டை
முழுமைப் பண்பு வெளிப்படுத்துகிறது என்று கண்டது ஓர் ஆழமான உட்பார்வை.
மேலும், மெய்யெண்களை டெடெகிண்ட் வெட்டுகளாக வெளிப்படையாகக்
கட்டமைப்பது, பகுப்பாய்வியலின் பாடப்பொருளுக்கு உறுதியான அடித்தளத்தைத்
தருகிறது. \emph{முழுமையான வரிசைப்படுத்தப்பட்ட புலம்} என்ற கருத்து
முரணற்றது என்பதை அறிவோம்; ஏனெனில் வெட்டுகள் அத்தகைய ஒரு புலமாக அமைகின்றன.

% SEGMENT OLP-0046-S02
இவை அனைத்தையும் ஏற்றாலும், இந்தச் சாதனைகளைப் பற்றி நமக்குள்ள
சில ஐயங்களையும் வெளிப்படுத்த வேண்டும்.

% SEGMENT OLP-0046-S03
மெய்யெண்களை அவற்றின் பழக்கமான (முடிவுறாததாக இருக்கக்கூடிய)
தசம விரிவாக்கங்களாகக் கருதுவதைவிட, வெட்டுகளாகக் கருதுவது
\emph{மேலும்} கடுமையானதா என்பது தெளிவில்லை.
மெய்யெண்களுக்கான இந்தப் பிந்தைய ``கட்டமைப்பு'',
கோஷி தொடர் மூலம் செய்யப்படும் மெய்யெண் கட்டமைப்புடன் சில ஒற்றுமைகளைக்
கொண்டுள்ளது; ஆனால் உண்மையில், அது 17ஆம் நூற்றாண்டின் தொடக்கத்திலிருந்தே
கணிதவியலாளர்களுக்குச் சாரத்தில் தெரிந்திருந்தது
(\olref[cauchy]{sec} ஐப் பார்க்கவும்).
மெய்யெண்களுக்கு முழுமைப் பண்பு உள்ளது என்பதை உணர்ந்ததிலிருந்தே
கடுமையில் உண்மையான உயர்வு வந்தது;
குறிப்பிட்ட கணங்களாக மெய்யெண்களைக் கட்டமைக்கும் திறன் மட்டும்,
அவ்வளவு சுவையானதாக இல்லாமல் இருக்கலாம்.

% SEGMENT OLP-0046-S04
முழுக்களையோ விகிதமுறு எண்களையோ (மிக எளிதாகச் செய்யப்படும்)
எண்முறைமைக் கட்டமைப்பு, அந்தப் பகுதிகளில் கடுமையை அதிகரிக்கிறது
என்பது இன்னும் தெளிவற்றது.
இங்கு ஓர் எளிய கவனிப்பைக் கூறுவது பயனுள்ளது.
முழுக்களை இயல் எண்களின் வரிசைப்படுத்தப்பட்ட ஜோடிகளின் சமான
வகுப்புகளாக \emph{கட்டமைத்து}, பின்னர் விகிதமுறு எண்களை முழு ஜோடிகளின்
சமான வகுப்புகளாகக் கட்டமைத்து, அதன்பின் மெய்யெண்களை விகிதமுறு எண்களின்
கணங்களாகக் கட்டமைத்தவுடன், அந்தக் கட்டமைப்புகளை உடனே
\emph{மறந்துவிடுகிறோம்}.
குறிப்பாக, அந்தக் கட்டமைப்புகள் செயல்பட வேண்டியவாறு செயல்படுகின்றன
என்பதை நிறுவுவதைத் தவிர, ஒரு கணித நிறுவலின்போது யாரும் அவற்றை
\emph{பயன்படுத்த} விரும்பமாட்டார்கள்.
இயல் எண்களின் கணங்களின் கணங்களின் கணங்களின் கணங்களின் கணங்களின்
கணங்களின் ஏதோ ஒரு கணத்தைப் பற்றிப் பேசுவதைவிட,
ஒரு மெய்யெண்ணைப் பற்றி நேரடியாகப் பேசுவது மிகவும் எளிது.
%(மெய்யெண்கள் விகிதமுறு எண்களின் கணங்களாகக் கட்டமைக்கப்பட்டன; அவை முழு ஜோடிகளின் சமான வகுப்புகளாக, அதாவது முழுக்களின் கணங்களின் கணங்களின் கணங்களாகக் கட்டமைக்கப்பட்டன; இவ்வாறே தொடர்கிறது.)
%2/3 = [2,3]
%	அதாவது {<2,3>, ... }
%	{ {{2},{2,3}}, ... }
%
%மெய்யெண்கள்
%	விகிதமுறு எண்களின் கணங்கள்
%
%விகிதமுறு எண்கள்
%	முழுக்களின் கணங்களின் கணங்களின் கணங்கள்
%
%முழுக்கள்
%	இயல் எண்களின் கணங்களின் கணங்களின் கணங்கள்
	
% SEGMENT OLP-0046-S05
இந்த வரையறைகள் முழுக்கள், விகிதமுறு எண்கள் அல்லது மெய்யெண்கள்
\emph{மெய்ப்பொருளியல் நோக்கில்} \emph{எவை} என்பதைச் சொல்கின்றனவா என்பதே
எல்லாவற்றையும்விட அதிக ஐயத்திற்குரியது.
அதாவது, மெய்யெண்கள் (எடுத்துக்காட்டாக) குறிப்பிட்ட
(கணங்களின் கணங்களின் \ldots) கணங்களாகவே \emph{உள்ளன} என்பது ஐயத்திற்குரியது.
இந்தக் கருத்திற்கான முக்கியத் தடை, கட்டமைப்பை வெவ்வேறு பல வழிகளில்
செய்திருக்க முடியும் என்பதாகும்.
மெய்யெண்களைப் பொறுத்தவரை, உண்மையாகவே சுவையான பல வேறுபட்ட
கட்டமைப்புகள் உள்ளன (\olref[cauchy]{sec} ஐப் பார்க்கவும்).
ஆனால் வேறுபட்ட கட்டமைப்புகளைப் பெறுவதற்கான மிக அற்பமான வழி இதோ:
\olref[sfr][rel][ref]{sec} இல் செய்ததுபோல,
வரிசைப்படுத்தப்பட்ட ஜோடிகளைச் சற்றே வேறுவிதமாக வரையறுத்திருக்கலாம்.
வரிசைப்படுத்தப்பட்ட ஜோடிக்கான இந்த மாற்றுக் கருத்தைப் பயன்படுத்தியிருந்தால்,
நமது கட்டமைப்புகள் முன்புபோலவே மிகச்சரியாகச் செயல்பட்டிருக்கும்;
ஆனால் இறுதியில் வேறு பொருள்கள் கிடைத்திருக்கும்.
எனவே முழுக்கள், விகிதமுறு எண்கள், மெய்யெண்கள் ஆகியவற்றுக்குப்
போட்டியான பல கணக்கோட்பாட்டுக் கட்டமைப்புகள் உள்ளன.
இப்போது முழுக்கள் (எடுத்துக்காட்டாக) \emph{இந்தக்} கணங்கள்தான்,
\emph{அந்தக்} கணங்கள் அல்ல என்று கூறுவது தன்னிச்சையானதாகவும்
தர்மசங்கடமூட்டுவதாகவும் இருக்கும்.
(\olref[sfr][rel][ref]{sec} இல் கண்டதுபோல,
இது \citealt{Benacerraf1965} புகழ்பெறச் செய்த ஒரு வாதத்தின் எடுத்துக்காட்டு.)

% SEGMENT OLP-0046-S06
மேலும் ஒரு கருத்தை எழுப்புவது பயனுள்ளது:
நமது கட்டமைப்புகளில் முற்றிலும் \emph{விநோதமான} ஒன்று உள்ளது.
இயல் எண்களிலிருந்து தொடங்கினோம்.
பின்னர் முழுக்களையும், ``முழுக்களின் $0$'' என்பதையும்,
அதாவது $ \equivrep{0,0}{\Intequiv}$ ஐயும் கட்டமைக்கிறோம்.
ஆனால் $0 \neq \equivrep{0,0}{\Intequiv}$.
உண்மையில், நமது கட்டமைப்புகளின்படி \emph{எந்த} இயல் எண்ணும் முழு அல்ல.
ஆனால் இது உள்ளுணர்விற்கு முற்றிலும் எதிரானதாகத் தோன்றுகிறது.
மேலும், \olref[sfr][set][imp]{sec} இல்,
$\Nat \subseteq \Rat$ என்று அதிக வாதமின்றிக் கூறினோம்.
எண்கள் மிகச்சரியாக \emph{எவை} என்பதை அந்தக் கட்டமைப்புகள் சொல்கின்றன
எனில், இந்தக் கூற்று அற்பமாகவே பொய்யானது.

% SEGMENT OLP-0046-S07
சற்று விலகி நின்று பார்த்தால், நாம் எங்கு வந்துள்ளோம்?
முறைசாராத கணக்கோட்பாட்டில் வேலைசெய்து,
இயல் எண்களை நமக்குத் தரப்பட்டவையாகப் பயன்படுத்தி,
முழுக்கள், விகிதமுறு எண்கள், மெய்யெண்கள் ஆகியவற்றைக் குறிப்பிட்ட
கணங்களாகக் \emph{கருத} முடிகிறது.
அந்தப் பொருளில், இந்தப் பொருள்களின் கோட்பாடுகளை ஒரு
கணக்கோட்பாட்டுக்குள் \emph{பதிக்க} முடிகிறது.
ஆனால் இந்தப் பதித்தலின் மெய்யியல் பொருள் அவ்வளவு நேரடியானதல்ல.

% SEGMENT OLP-0046-S08
நிச்சயமாக, இவை எதுவும் இறுதி முடிவல்ல!{}
கருத்து இதுமட்டுமே.
மெய்யெண்களின் எண்முறைமைக் கட்டமைப்பு ஆழமான மெய்யியல் முக்கியத்துவம்
\emph{கொண்டுள்ளது} எனக் காட்டுவதற்கு,
மேலும் சில \emph{மெய்யியல்} வாதங்கள் தேவைப்படும்.

%மேலும்: இந்த அத்தியாயம் முழுவதிலும், இயல் எண்கள் நமக்கு வெறுமனே
%\emph{தரப்பட்டவை} எனக் கருதியுள்ளோம். ஆனால் அவற்றை வைத்து என்ன செய்யலாம்?
%அதுவே அடுத்த அத்தியாயத்தின் பாடப்பொருள்.

\end{document}
